%==============================================================
% INTRODUCTION + RELATED WORK -- journal version.
%
% Rewritten against the conference version because the contributions changed:
% the benchmark decomposition and the artefact control are new, and two of the
% conference version's own conclusions are corrected here.
%
% THREE NEW REFERENCES, each verified against a machine-readable record:
%   shortcut     Geirhos et al., Nature Machine Intelligence 2020,
%                doi:10.1038/s42256-020-00257-z  (arXiv:2004.07780)
%   datasetbias  Torralba & Efros, CVPR 2011, pp. 1521-1528,
%                doi:10.1109/CVPR.2011.5995347
%   nadeau       Nadeau & Bengio, Machine Learning 52(3), 2003, pp. 239-281,
%                doi:10.1023/A:1024068626366
% Bibitems for all three are at the end of this file.
%==============================================================

\section{Introduction}

Bangladesh's social media ecosystem has become a primary vector for manipulated
news imagery, from digitally altered photocards and forged documents to
AI-generated video frames misattributed to public figures. Unlike text-only
misinformation, this content exploits the visual and textual channels jointly:
a viral image is judged not on its pixels alone but on the claim attached to it
as it spreads.

The problem is compounded by linguistic specificity. Most fake-news detection
methods are built on English datasets from Western platforms and presuppose
signals---dense comment threads, public retweet graphs---that are sparse or
structured differently in Bangladesh. No large-scale, organically collected
corpus of manipulated Bangladeshi news images exists, so training must rely on
synthetically constructed benchmarks.

That reliance is the starting point of this paper. A constructed benchmark is
reproducible, which is its virtue, but it is also assembled by a pipeline, and
a pipeline leaves traces. Two consequences follow, and neither is checked as a
matter of course. First, the classes may be separable for reasons that have
nothing to do with the phenomenon under study: if the generator emits images at
one resolution and the authentic images arrive at many, a model can score
perfectly without looking at a picture---the shortcut-learning failure
\cite{shortcut} in the form dataset-bias work has warned about for over a
decade \cite{datasetbias}, and which audits of fake-news corpora have since
found in practice \cite{verma}. Second, a benchmark that advertises several
``attack types'' may not pose the same task in each of them: if an attack
rewrites an article's text and reuses its photograph unchanged, then an
image-only model is being asked a question whose answer is not present in its
input, and whatever number it produces is not a measure of its competence. A
related concern, unimodal bias, has been documented in English multimodal
misinformation benchmarks \cite{verite}.

We find both in BanFakeMM \cite{banfakemm}, the benchmark this work builds on,
and both change what the surrounding results mean. This paper therefore reports
a detector and, inseparably, the measurement discipline needed to say anything
about it.

Our contributions are as follows.

\begin{itemize}
\item \textbf{A decomposition of the benchmark by modality.} Of BanFakeMM's
five adversarial splits, two manipulate the image and three manipulate the
article text while reusing the photograph: in each of those three, every one
of the 1{,}000 manipulated images is byte-identical to an authentic image from
the benchmark's own corpus. An image-only detector is well posed on two splits
and undecidable on three, and we show its measured behaviour matches that
prediction to the decimal.

\item \textbf{A dataset artefact, and a control that it survives.} In the
training split every manipulated image shares one resolution, and a classifier
given only file metadata---no pixels---separates the classes perfectly.
Re-encoding every image through one identical pipeline removes the artefact and
costs the visual model $1.03$ accuracy points, which establishes that the
signal is in the content.

\item \textbf{A per-modality ablation under duplicate-aware grouped
cross-validation}, which removes OCR on evidence rather than by assumption and
shows the caption-fusion pathway does not earn its computational cost, together
with a like-for-like comparison against the feature families of the closest
prior Bangla multimodal system under that same protocol.

\item \textbf{A distribution-shift evaluation} on a 1{,}308-image sealed
partition---$6.7\times$ the one used in earlier work on this system---with the
artefact controlled, thresholds selected on a disjoint partition, and
article-level overlap measured rather than assumed away.

\item \textbf{Evidence that supervised adaptation helps}, and an account of why
an earlier measurement said otherwise. Adaptation on 300 target-style images
improves ranking by $+0.025$ AUROC with confidence intervals excluding zero
across eight seeds, in both the controlled and uncontrolled settings. The
earlier null rested on accuracy at a fixed threshold on 194 images from a
single fit---a statistic we show can disagree in sign with the effect it is
meant to measure.
\end{itemize}

Two of these findings correct conclusions we previously published about this
same system. We state them as corrections, with the measurements that force
them, rather than quietly revising the numbers.

We claim no architectural novelty. BLIP/CLIP feature extraction with a
gradient-boosted classifier is well established, and the accuracy it reaches
here is unremarkable. What we have not found in the Bangla literature is an
image-level manipulation detector whose reported numbers are accompanied by
evidence that the benchmark supports them.

\section{Related Work}

\subsection{Bangla Fake News Detection}

Research on Bangla-language fake news detection remains limited relative to
English. BanFakeNews \cite{banfakenews} established an early large-scale
benchmark of roughly 50{,}000 Bangla news items labeled authentic or
fabricated, and BanFakeNews-2.0 \cite{banfakenews2} extended that line to a
60{,}000-instance curated corpus with an LLM-based benchmark. Both are
text-only. MultiBanFakeDetect \cite{multibanfake} moved to multimodal
settings with a hybrid DenseNet-169/mBERT fusion architecture and is the
closest prior system to the present work in both language and modality scope.
Beyond Bangla specifically, MMCFND and its MMIFND dataset \cite{mmcfnd} extend
caption-aware detection to seven low-resource Indic languages including
Bengali.

None of these is organised around a defined image-manipulation task with
adversarially constructed positive examples: MultiBanFakeDetect labels whole
articles by their overall claim, so a model may succeed on textual cues alone,
and the remaining systems do not take images as a decision input at all. This
is the gap BanFakeMM \cite{banfakemm} addresses and the one we work in---with
the qualification, established in Section~\ref{sec:detectable}, that only part
of that benchmark poses an image-level task.

\subsection{Multimodal Misinformation Detection}

Outside the Bangla literature, EANN \cite{eann}, MVAE \cite{mvae} and SpotFake
\cite{spotfake} established that visual and textual signals are complementary
for misinformation detection, each on English-language, Western-platform data.
A more recent line pairs multimodal classification with retrieval-augmented,
multi-model verification---DEFAME \cite{defame}, RAGAR \cite{ragar}, RAMA
\cite{rama} and CaLM \cite{calm}---and reports gains from grounding a verdict
in retrieved evidence rather than in the classifier alone. None is
Bangla-specific or grounded in a Bangladeshi corpus. We do not evaluate a
verification stage here; Section~\ref{sec:limitations} explains what
evaluating one properly would require and why we judge our preliminary
measurement insufficient to report.

Work on adapting vision-language models under distribution shift is directly
relevant to our shift results. CLIPood \cite{clipood} adapts CLIP itself for
out-of-distribution generalization, and DPOD \cite{dpod} learns
domain-specific prompts for multimodal misinformation detection using
out-of-domain data. Both are architecture-level interventions; our shift
mitigation is data-centric---threshold selection on a disjoint partition, and
supervised adaptation on target-style examples---and we do not claim the two
are exclusive.

\subsection{Evaluation Integrity in Constructed Benchmarks}
\label{sec:related-integrity}

The concerns that organise this paper are not new to machine learning, only
under-applied in this subfield. Torralba and Efros \cite{datasetbias} showed
that visual datasets carry signatures strong enough for a classifier to name
the dataset an image came from, and argued that benchmark performance can
therefore measure dataset membership rather than the target concept. Geirhos
et al.\ \cite{shortcut} generalised this as shortcut learning: networks
reliably prefer any decision rule that separates the training classes, however
incidental, and such rules are invisible to accuracy on an
identically-constructed test set. The diagnostic they recommend---evaluate on
data where the shortcut is absent---is what our normalisation control
(Section~\ref{sec:normalise}) implements, and what our per-split decomposition
(Section~\ref{sec:detectable}) does at the level of the task rather than the
features.

The same concerns have been raised inside misinformation detection, though not
for Bangla and not for image benchmarks built this way. Verma \cite{verma}
audits the ISOT fake-news corpus and finds that a classifier given only the
subject field, with the article text removed, reaches an F1 of $1.000$ because
the two classes have disjoint subjects; removing the identified leakage
channels barely moves F1, but a topic-disjoint split lowers it substantially
and on an independent benchmark all models fall to near-chance ranking. The
paper recommends metadata-only and topic-disjoint baselines as cheap
diagnostics, which is the diagnostic we apply to image files. For multimodal
benchmarks, Papadopoulos et al.\ \cite{verite} showed that two widely used
ones, VMU-Twitter and COSMOS, carry unimodal bias---a detector reading one
modality can match or beat multimodal ones---and built VERITE so that each
image and each caption occurs in both a truthful and a misleading pair; Lin et
al.\ \cite{linmodality} propose automated analyses that quantify modality bias
per sample. These approaches measure, through trained detectors, how much each
modality contributes. Our per-split analysis is complementary: where an attack
changes only the article text, the answer for an image-only model follows
exactly from content hashes, with nothing trained.

The statistical half of the problem is equally well established and equally
often skipped. Nadeau and Bengio \cite{nadeau} showed that the folds of a
resampled cross-validation are not independent, because their training sets
overlap, and that the naive variance estimate is consequently too small; they
give the correction we apply in Section~\ref{sec:training}. Our results include
a concrete instance of the cost of ignoring this: a null result reported on a
small test partition, at a single threshold, from a single fit, which a
properly powered and threshold-free measurement reverses
(Section~\ref{sec:adaptation}).

To our knowledge no prior Bangla multimodal detection work reports a
metadata-only control, a per-split analysis of which modality carries the
manipulation, or a correction for fold dependence. The ideas are not new:
metadata-only probes and unimodal-bias analysis exist for English text and
social-media corpora \cite{verma,verite}. What we add is their application to a
Bangla image-manipulation benchmark, and an exact test---byte-level reuse of
the authentic image---for which attack types an image-only model can in
principle detect. We take these to be the minimum reporting standard for
results on constructed benchmarks of this kind, rather than contributions
peculiar to this paper.

%==============================================================
% BIBITEMS for the three new references. Insert into the bibliography;
% IEEE numbering follows order of first citation.
%==============================================================
%
% \bibitem{shortcut} R. Geirhos, J.-H. Jacobsen, C. Michaelis, R. Zemel,
% W. Brendel, M. Bethge, and F. A. Wichmann, ``Shortcut learning in deep neural
% networks,'' \emph{Nature Machine Intelligence}, vol. 2, pp. 665--673, 2020,
% doi: 10.1038/s42256-020-00257-z.
%
% \bibitem{datasetbias} A. Torralba and A. A. Efros, ``Unbiased look at dataset
% bias,'' in \emph{Proc. IEEE Conf. Computer Vision and Pattern Recognition
% (CVPR)}, 2011, pp. 1521--1528, doi: 10.1109/CVPR.2011.5995347.
%
% \bibitem{nadeau} C. Nadeau and Y. Bengio, ``Inference for the generalization
% error,'' \emph{Machine Learning}, vol. 52, no. 3, pp. 239--281, 2003,
% doi: 10.1023/A:1024068626366.
